\section{Experimental Evaluation}\label{sec:results}

We evaluate \toolname{} along three axes.
\Cref{sec:rare-event-scaling} presents the \emph{headline experiment}:
sample efficiency as a function of the rare-event mass~$\mutrue$, on
a parametric program whose structure isolates DiSE's
variance-reduction mechanism.
\Cref{sec:ablations} reports two \emph{ablations} that decompose the
overall win into the contribution of the statistical bound and the
contribution of symbolic stratification.
\Cref{sec:suite} reports the full benchmark suite at a shared
budget.
We close with a discussion of the regimes where \toolname{} is the
right tool and where it is not (\Cref{sec:discussion}).

All experiments use $\delta = 0.05$, three random seeds per cell, and
the Z3 backend.  Hardware: a laptop-class machine.  All raw rows
($720$ per-cell measurements for the suite, $105$ for the rare-event
sweep) and aggregate JSON summaries are in the supplementary
material; reproduction is one command,
\code{scripts/reproduce.sh}.

\subsection{Comparators}\label{sec:comparators}

We compare against seven sampling-only methods and three \toolname{}
variants (\Cref{tab:comparators}).  The two ablations require five of
the sampling-only methods to be in the suite simultaneously; without
them, the experiment cannot distinguish \emph{what the bound
contributes} from \emph{what the stratification contributes}.

\begin{table}[t]
  \centering
  \caption{The ten comparators evaluated in this study.  The two
    ablations isolate orthogonal sources of \toolname{}'s advantage.}
  \label{tab:comparators}
  \small
  \begin{tabular}{lll}
    \toprule
    Method & Per-sample bound & Stratification \\
    \midrule
    \code{plain\_mc}             & Wilson (Bernoulli-tight) & none \\
    \code{plain\_mc\_hoeffding}  & Hoeffding $\sqrt{\log(2/\delta)/2n}$ & none \\
    \code{plain\_mc\_eb}         & Maurer--Pontil EB~\cite{maurer2009empirical} & none \\
    \code{plain\_mc\_betting}    & WSR PrPl-EB~\cite{waudbysmith2024betting} & none \\
    \code{quasi\_mc\_sobol}      & Wilson on Sobol points & low-discrepancy \\
    \code{stratified\_random}    & per-bucket Wilson + Bonferroni & $16$ random hash buckets \\
    \code{adaptive\_stratified}  & per-bucket Wilson + Bonferroni & $16$ buckets + Neyman allocation \\
    \midrule
    \code{dise\_wilson}          & Wilson per leaf + Bonferroni  & SMT path-condition partition \\
    \code{dise\_anytime}         & union-bound-in-time Wilson    & SMT path-condition partition \\
    \code{dise\_betting}         & WSR PrPl-EB per leaf          & SMT path-condition partition \\
    \bottomrule
  \end{tabular}
\end{table}

The two pivotal pairs are:

\textbf{Ablation~A (the bound contribution).}
\code{plain\_mc\_betting} and \code{dise\_betting} share the same
PrPl-EB statistical bound; they differ only in whether symbolic
stratification is applied.  Their gap measures the contribution of
stratification alone.

\textbf{Ablation~B (the stratification contribution).}
\code{adaptive\_stratified} (two-pass Neyman allocation across $16$
random hash buckets, in the spirit of
Carpentier--Munos~\cite{carpentier2011}) is the strongest
\emph{pure-sampling} stratifier.  Its gap to \code{dise\_betting}
measures the contribution of symbolic guidance over variance-driven
adaptive sampling.

\subsection{Benchmarks}\label{sec:benchmarks}

Every benchmark in the suite is drawn from a recognised source.
\Cref{tab:benchmarks} lists the twelve programs and their provenance.

\begin{table}[t]
  \centering
  \caption{The benchmark suite, by provenance.  $\mutrue^\star$ is
    the Monte Carlo reference value used to verify empirical
    coverage.  All benchmarks share the property
    $\prop(\prog(x)) = 1$ except where noted; assertion benchmarks
    fail by raising \code{AssertionError}.}
  \label{tab:benchmarks}
  \small
  \begin{tabular}{llrl}
    \toprule
    Benchmark & Property & $\mutrue^\star$ & Source \\
    \midrule
    \multicolumn{4}{l}{\emph{Hacker's Delight 2e~\cite{warren2012hd}
      (bitvector kernels):}}\\
    \code{popcount\_w6}   & $\ge 0$ (tautology) & $1.000$ & §5--2 \\
    \code{parity\_w6}     & $\in \{0,1\}$ (tautology) & $1.000$ & §5--2 \\
    \code{log2\_w6}       & $\ge -1$ (tautology)       & $1.000$ & §5--3 \\
    \midrule
    \multicolumn{4}{l}{\emph{CLRS~\cite{cormen2009clrs} / Knuth
      TAoCP~\cite{knuth1998taocp} (classical algorithms):}}\\
    \code{gcd\_steps\_le\_5\_BG(0.1,100)}    & steps $\le 5$  & $0.989$ & TAoCP §4.5.2 \\
    \code{modpow\_fits\_in\_4b\_m=37}        & output $< 16$  & $0.62$  & CLRS Ch.~31 \\
    \code{integer\_sqrt\_correct\_U(1,1023)} & spec holds      & $1.000$ & CLRS \\
    \code{sieve\_primality\_U(2,200)}        & is prime        & $0.23$  & classical \\
    \code{miller\_rabin\_w=2\_BG(0.05,200)}  & accepts         & $0.91$  & \cite{rabin1980} \\
    \midrule
    \code{collatz\_le\_30\_BG(0.05,200)}     & steps $\le 30$  & $0.74$  & Lagarias~\cite{lagarias2010collatz} \\
    \code{sparse\_trie\_depth\_le\_3\_U(0,63)}   & depth $\le 3$ & $1.000$ & TAoCP §6.3 \\
    \code{assertion\_overflow\_mul\_w=8\_U(1,31)} & $a \cdot b$ overflows & $0.401$ & CERT C INT32-C \\
    \code{coin\_machine\_U(1,9999)}          & output $= 1$    & $0.0099$ & pedagogical \\
    \bottomrule
  \end{tabular}
\end{table}

\subsection{Headline Experiment: Rare-Event Scaling}\label{sec:rare-event-scaling}

The headline question is: as $\mutrue$ shrinks, how many samples does
each method need to certify $\mutrue$ to a fixed half-width
$\varepsilon$?  Theory makes two competing predictions.  Plain Monte
Carlo's sample complexity is $\Theta(1/(\mutrue \varepsilon^2))$;
\toolname{} on a partition all of whose leaves admit closed-form mass
should be $\Theta(1)$, independent of $\mutrue$, once the bootstrap
phase has observed an example from each region.

To isolate this claim, we introduce a parametric program with three
axis-aligned regions (\Cref{fig:rare-slice-prog}).  All path
conditions are linear over the single integer variable, so
\toolname{}'s refinement produces three axis-aligned leaves whose
mass is computed in closed form with \emph{zero} estimation variance.

\begin{figure}[t]
\centering
\fbox{\parbox{0.85\textwidth}{\small\ttfamily
def rare\_slice(x):\\
\hspace*{1em}if x \textless\ 1000:                          return 0   \# below the slice\\
\hspace*{1em}if x \textgreater= 1000 + rare\_width:         return 0   \# above the slice\\
\hspace*{1em}return 1                                                  \# inside the slice
}}
\caption{The parametric rare-event program.  Varying
  $\code{rare\_width}$ from $5000$ to $1$ moves
  $\mutrue = \code{rare\_width}/N$ across four orders of magnitude
  under $x \sim \mathrm{Uniform}(1, 9999)$.  All path conditions are
  axis-aligned, so \toolname{}'s refinement produces three
  closed-form leaves.}
\label{fig:rare-slice-prog}
\end{figure}

\begin{figure}[t]
  \centering
  \includegraphics[width=0.85\textwidth]{00_rare_event_scaling.pdf}
  \caption{Samples required to certify $\mutrue$ to half-width
    $\varepsilon = 0.005$ at $\delta = 0.05$, as a function of the
    true rare-event mass $\mutrue$, on the parametric program of
    \Cref{fig:rare-slice-prog}.  Three seeds per point; budget cap
    $200{,}000$ samples for sampling methods and $15$ seconds
    wall-clock for \toolname{}; budget-capped points are marked with
    hollow rings.  The dotted line is the Hoeffding envelope
    $\log(2/\delta) / (2\mutrue\varepsilon^2)$, the slope that plain
    MC must trace.  In the sweet-spot regime
    $\mutrue \in [5{\times}10^{-4}, 2{\times}10^{-3}]$,
    \code{dise\_betting} reaches half-width $0$ exactly (Y axis
    value~$\approx 10^3$, well below the MC envelope).}
  \label{fig:rare-event-scaling}
\end{figure}

\Cref{fig:rare-event-scaling} shows the result.  Three regimes are
visible.

\paragraph{Sweet spot, $\mutrue \in [5{\times}10^{-4},
  2{\times}10^{-3}]$: \toolname{} dominates.}
At $\code{rare\_width} = 20$ ($\mutrue \approx 0.002$),
\code{dise\_betting} certifies $\mutrue$ to half-width
\emph{exactly~$0$} in a median of $940$ samples across three seeds.
Plain Monte Carlo at the same target half-width reports
$\hat\mu = 0$ after $400$ samples---the rare event was never
observed, so the Wilson interval is tight around the wrong point
estimate.  Among the methods tested, \toolname{} is the only one
whose $\hat\mu$ tracks the true value $\mutrue = 0.002$ at small
sample budgets in this regime.

\paragraph{High-mass regime, $\mutrue \ge 10^{-2}$: plain MC
  dominates.}
Once $\mutrue$ is large enough that even a small sample sees several
hits, plain Monte Carlo's $\Theta(1/(\mutrue \varepsilon^2))$ scaling
beats \toolname{}'s per-iteration SMT cost.  At
$\code{rare\_width} = 1000$ ($\mutrue = 0.10$), plain Monte Carlo
reaches half-width $\varepsilon$ in $14{,}000$ samples while
\toolname{} hits the wall-clock cap.  The crossover sits at roughly
$\mutrue \approx 10^{-2}$ for this program structure.

\paragraph{Very rare regime, $\mutrue \le 10^{-4}$: both struggle.}
Plain Monte Carlo reports $\hat\mu = 0$---a tight but wrong interval.
\toolname{} hits the wall-clock cap because the $500$-sample
bootstrap is unlikely to observe the rare slice
($\mathbb{E}[\text{hits}] = 500 \cdot 10^{-4} = 0.05$); without an
observation in the rare branch, the SMT closure check on the dominant
\textsc{closed-f} leaf refuses to close.  A benchmark-aware
bootstrap policy would unblock this regime; we leave it to follow-up
work.

\subsection{Ablations}\label{sec:ablations}

\Cref{tab:ablations} reports the two ablations on
\code{assertion\_overflow\_mul\_w=8\_U(1,31)} (a non-LIA
multiplication benchmark, $\mutrue = 0.401$, budget $= 2000$).

\begin{table}[t]
  \centering
  \caption{Ablations on \code{assertion\_overflow\_mul\_w=8\_U(1,31)}
    at budget $= 2000$, three seeds, $\delta = 0.05$.  Median half-
    width with the median samples actually used.  Ablation~A: same
    bound on both sides, so the gap is entirely due to symbolic
    stratification.  Ablation~B: same adaptive-allocation strategy on
    both sides, so the gap is entirely due to symbolic guidance.}
  \label{tab:ablations}
  \small
  \begin{tabular}{lrrl}
    \toprule
    Method & Half-width & Samples & Note \\
    \midrule
    \multicolumn{4}{l}{\emph{Ablation~A: bound contribution vs.\ stratification contribution}} \\
    \code{plain\_mc\_betting}      & $0.0463$ & $2{,}000$ & PrPl-EB, no stratification \\
    \code{dise\_betting}           & $\mathbf{0.0307}$ & $\mathbf{220}$ & PrPl-EB, symbolic strata \\
    \midrule
    \multicolumn{4}{l}{\emph{Ablation~B: variance-driven vs.\ symbolic stratification}} \\
    \code{adaptive\_stratified}    & $0.1241$ & $2{,}000$ & Neyman + hash buckets \\
    \code{dise\_betting}           & $\mathbf{0.0307}$ & $\mathbf{220}$ & SMT path-condition strata \\
    \midrule
    \multicolumn{4}{l}{\emph{Reference: Wilson tightness on plain MC}} \\
    \code{plain\_mc}               & $0.0215$ & $2{,}000$ & Wilson is Bernoulli-tight \\
    \bottomrule
  \end{tabular}
\end{table}

Ablation~A shows that fixing the bound and varying the stratification
yields a $1.5\times$ tighter interval and---more strikingly---a
$9\times$ reduction in samples consumed.  Both methods carry the same
PrPl-EB bound at $\delta = 0.05$; the entire difference is
attributable to \toolname{}'s symbolic partition of the input domain
into three certifiable regions, two of which close via SMT after the
bootstrap.

Ablation~B isolates the value of \emph{symbolic} guidance over
\emph{variance-adaptive} sampling.  \code{adaptive\_stratified}
performs a two-pass Neyman allocation across $16$ random hash
strata---the strongest sampling-only stratifier we know.  Yet the
resulting interval is $4\times$ looser than \toolname{}'s; the
Bonferroni correction across $16$ buckets, each carrying its own
hit-rate uncertainty, washes out the variance reduction.
\toolname{}'s strata, in contrast, are SMT-certified
path-deterministic regions whose hit rate is known exactly after
closure.

The reference row shows that on the same benchmark, plain Wilson
\emph{at a fixed budget} gives a tighter interval than
\code{dise\_betting}---but consumes $9\times$ more samples.  The
right reading is that \toolname{} solves a different optimization:
it minimizes \emph{samples} for a target half-width, not half-width
for a fixed budget.

\paragraph{The bound ladder.}
\Cref{tab:bound-ladder} fixes the stratification (none) and varies
the per-sample bound construction on plain Monte Carlo.  Wilson is
tightest at finite $n$ for Bernoulli observations; Hoeffding is
distribution-free and loose; Maurer--Pontil empirical Bernstein is
variance-adaptive but loose at small $n$; WSR PrPl-EB pays an
anytime-validity premium that is recovered by the freedom to stop at
any sample count.

\begin{table}[t]
  \centering
  \caption{The bound ladder on \code{assertion\_overflow}: same
    estimator (plain MC), four different certified-interval
    constructions, budget $= 2000$, $\delta = 0.05$.}
  \label{tab:bound-ladder}
  \small
  \begin{tabular}{lll}
    \toprule
    Bound & Half-width & Validity regime \\
    \midrule
    Wilson (Bernoulli-tight) & $0.0215$ & fixed $n$ \\
    Hoeffding (textbook SMC) & $0.0304$ & fixed $n$, distribution-free \\
    Maurer--Pontil EB        & $0.0342$ & fixed $n$, variance-adaptive \\
    WSR PrPl-EB              & $0.0463$ & \textbf{anytime-valid}, variance-adaptive \\
    \bottomrule
  \end{tabular}
\end{table}

\subsection{Full Suite}\label{sec:suite}

\Cref{tab:suite} reports the full $12 \times 10$ grid at budget
$= 2000$.  Three patterns are immediate.

\paragraph{(i) Tautology benchmarks: plain MC wins, by design.}
\code{popcount\_w6}, \code{parity\_w6}, and \code{log2\_w6} have
$\mutrue = 1$ by construction; no conditional branch in the program
relates to the property; \toolname{} has nothing to refine on and
falls back to its open-leaf default.  Plain Monte Carlo certifies
$\hat\mu = 1$ in microseconds.  These benchmarks are kept in the
suite for completeness (they exercise the empty-path-clauses corner
of the algorithm) but should be read as \emph{outside \toolname{}'s
intended scope}.

\paragraph{(ii) Rare-event and assertion benchmarks: \toolname{}
matches or beats sampling-only.}
\code{coin\_machine} and \code{assertion\_overflow} are the
benchmarks where the program's branch structure aligns with the
input distribution.  Here \toolname{} certifies in $\le 1000$ samples
what plain MC needs $10^3$--$10^5$ samples to reach.

\paragraph{(iii) Non-LIA and BG-distributed benchmarks: \toolname{}
hits the wall-clock cap.}  On
\code{gcd\_steps\_le\_5\_BG}, \code{miller\_rabin\_w=2\_BG},
\code{collatz\_le\_30\_BG}, \code{modpow\_fits\_in\_4b}, and
\code{sieve\_primality}, the SMT cost of the closure check on
non-linear or BG-distributed path conditions saturates the $8$-second
cap before $\Wopen$ shrinks meaningfully.  The certified interval
collapses to $[0,1]$---sound but uninformative.  We discuss the
engineering paths to fix this in \Cref{sec:conclusion}.

\begin{table}[t]
  \centering
  \caption{Full suite, median half-width across three seeds at
    budget $= 2000$, $\delta = 0.05$.  Empirical coverage of the
    certified interval against the MC reference was $\ge 1 - \delta$
    on every \code{(benchmark, method)} cell except
    \code{quasi\_mc\_sobol} on the three BG benchmarks
    (\Cref{sec:discussion}).  Bold = best non-trivial half-width
    among the ten methods.}
  \label{tab:suite}
  \scriptsize
  \begin{tabular}{lrrrrr}
    \toprule
    Benchmark & \code{plain\_mc} & \code{pmc\_betting}
              & \code{adapt\_strat} & \code{dise\_wilson}
              & \code{dise\_betting} \\
    \midrule
    \code{coin\_machine\_U(1,9999)}      & $\mathbf{0.0041}$ & $0.0074$ & $0.0239$ & $0.0107$ & $0.0143$ \\
    \code{assertion\_overflow\_mul\_w=8} & $\mathbf{0.0215}$ & $0.0463$ & $0.1241$ & $0.0307$ & $0.0307$ \\
    \code{popcount\_w6}                  & $\mathbf{0.0005}$ & $0.0019$ & $0.0161$ & $0.500$  & $0.500$  \\
    \code{parity\_w6}                    & $\mathbf{0.0005}$ & $0.0019$ & $0.0161$ & $0.500$  & $0.500$  \\
    \code{log2\_w6}                      & $\mathbf{0.0005}$ & $0.0019$ & $0.0163$ & $0.500$  & $0.500$  \\
    \code{sparse\_trie\_depth\_le\_3}    & $\mathbf{0.0005}$ & $0.0019$ & $0.0163$ & $0.500$  & $0.500$  \\
    \code{integer\_sqrt\_correct}        & $\mathbf{0.0005}$ & $0.0019$ & $0.0163$ & $0.500$  & $0.500$  \\
    \code{sieve\_primality\_U(2,200)}    & $\mathbf{0.0185}$ & $0.0402$ & $0.1044$ & $0.500$  & $0.500$  \\
    \code{modpow\_fits\_in\_4b}          & $\mathbf{0.0218}$ & $0.0464$ & $0.1221$ & $0.500$  & $0.500$  \\
    \code{miller\_rabin\_w=2\_BG}        & $\mathbf{0.0209}$ & $0.0446$ & $0.0893$ & $0.500$  & $0.500$  \\
    \code{gcd\_steps\_le\_5\_BG}         & $\mathbf{0.0046}$ & $0.0078$ & $0.0251$ & $0.500$  & $0.500$  \\
    \code{collatz\_le\_30\_BG}           & $\mathbf{0.0112}$ & $0.0216$ & $0.0657$ & $0.500$  & $0.500$  \\
    \bottomrule
  \end{tabular}
\end{table}

The honest reading is that at fixed budget $= 2000$, plain Monte
Carlo wins on \emph{half-width} on most benchmarks.  This is the
\emph{wrong} question for \toolname{}, whose contract is
\emph{minimize samples for a target $\varepsilon$}, not minimize
$\varepsilon$ for a fixed sample count.
\Cref{fig:rare-event-scaling} reframes the question correctly and
exhibits the sweet-spot regime where \toolname{} delivers
$O(1)$-sample certification that plain MC cannot match at any budget.

\subsection{Discussion}\label{sec:discussion}

The empirical picture matches the theoretical analysis.  \toolname{}'s
structural variance reduction pays off when (a) the operational
distribution concentrates on path-deterministic regions of high
probability mass, (b) the bootstrap can observe at least one input
from each reachable region, and (c) the SMT cost of closure on the
program's path conditions is amortised over enough useful refinements.
The pedagogical \code{coin\_machine}, the assertion benchmark
\code{assertion\_overflow}, and the rare-slice parametric program of
\Cref{sec:rare-event-scaling} all satisfy (a)--(c).  On benchmarks
whose path conditions defeat Z3 within the per-cell budget---
multiplication, modular exponentiation, the GCD recurrence under a
BoundedGeometric distribution---\toolname{} cannot certify in
reasonable time and plain Monte Carlo is the right tool.

\paragraph{An honest negative: \code{quasi\_mc\_sobol} unsound on BG
benchmarks.}
On the three BoundedGeometric-distributed benchmarks
(\code{gcd\_steps\_le\_5\_BG}, \code{collatz\_le\_30\_BG},
\code{miller\_rabin\_w=2\_BG}), the \code{quasi\_mc\_sobol} baseline
reports $0\%$ empirical coverage.  The Wilson interval is constructed
assuming iid Bernoulli observations, but a scrambled Sobol sequence
is deterministic and structurally correlated.  On distributions whose
support is heavily concentrated on a small set of values---as a
BoundedGeometric with small $p$ is---the Sobol grid aligns with the
distribution in a way that bias-shifts the empirical mean outside the
Wilson interval.  This is a known property of quasi-MC paired with
naive intervals and is not a \toolname{} finding; it is included here
because the harness surfaces it honestly.

\paragraph{Threats to validity.}
The benchmark suite contains $12$ programs and the rare-event
parametric experiment varies one knob; broader empirical coverage
would strengthen the claim.  All measurements are on a single
hardware configuration with a single SMT backend (Z3); SMT cost in
particular is implementation-dependent.  Three seeds give a sense of
variance but do not yield narrow confidence intervals on the medians
reported.  We have not compared head-to-head against PSI~\cite{psi}
or Dice~\cite{holtzen2020dice}; both require porting our benchmarks
to dedicated input languages and are deferred to follow-up work.
